%=============================================================================
% MÓDULO 01 — RESUMO COM FICHA CRÍTICA DAS FONTES AUDITÁVEIS
%=============================================================================
% Padrão: Cada resumo deve citar as duas fontes "ativas e auditáveis"
% fundamentais e apresentar fichas críticas resumidas.
%=============================================================================

% --- RESUMO DO ARTIGO ---
% Em português (obrigatório para submissão nacional) e inglês (opcional para internacional).

% Resumo em Português
\begin{abstract}
\textbf{Objetivo:} Investigar o conceito de assinatura cognitiva dentro de ecossistemas de sistemas orchestrados agentico, utilizando a arquitetura Core (MarceloClaro / OpenCode Ecosystem Core) como estudo de caso.

\textbf{Método:} Análise de dados técnicos do Evolution Registry (854 ciclos), Trust Engine scores, Spec Engine compliance e MetaBus episodic memory, fundamentado exclusivamente nas duas fontes "ativas e auditáveis": Clark & Chalmers (1998) e Verbveek (2005).

\textbf{Resultado Principal:} A assinatura cognitiva é reinterpretada como "rastreador de extensão protética" — padrão interpretativo documentável de como a cogniência humana se media, estende e se articula com uma orquestradora agente, e não como propriedade de mente ou personalidade do sistema.

\textbf{Implicações:} Descentralização da autoria para a parceria humano-prótese; accountability distribuído entre designer, usuário e agentes; sugestão de veículos de publicação alinhados a Filosofia da Tecnologia e HCI.
\end{abstract}

% Resumo em Inglês (para submissão internacional)
\begin{abstract}
\textbf{Objective:} Investigate the concept of cognitive signature within agentico orchestrated system ecosystems, using the Core architecture (MarceloClaro / OpenCode Ecosystem Core) as a case study.

\textbf{Method:} Analysis of technical data from Evolution Registry (854 cycles), Trust Engine scores, Spec Engine compliance and MetaBus episodic memory, funded exclusively on the two "active and auditables" primary sources: Clark \& Chalmers (1998) and Verbveek (2005).

\textbf{Main Result:} Cognitive signature is reinterpretated as "prosthetic extension tracker" — interpretable documented pattern of how human cognition is mediated, extended and articulated with an agente orchestrator, and not as mind or personality property of the system.

\textbf{Implications:} Decentralization of authorship to the human-prosthesis partnership; distributed accountability between designer, user and agents; publication venue suggestions aligned with Philosophy of Technology and HCI.
\end{abstract}

% Fichas Críticas Resumidas no Rodapé do Resumo
% Estas são as fichas críticas das duas fontes auditáveis, resumidas para o abstract
% Fonte 1 auditável: Clark & Chalmers (1998) - fundamenta a metáfora da prótese cognitiva
% Fonte 2 auditável: Verbveek (2005) - fundamenta a mediação ética e material da tecnologia
\end{abstract}